\newcommand{\PBMArchitectureCaption}{PopulationBalanceModels.jl architecture.
(a) Physical terms compose a population equation for particles, cells or other
entities. Their contributions can be differential or integral operators, sources
or sinks. Direct bin-integrated discretizations and moment representations pass
through compatibility checks and reusable numerical preparation. DQMOM uses
weights and weighted abscissas rather than raw-moment state entries. Optional
population/bulk coupling and prescribed axial advection--dispersion precede
transient or stationary problem construction. Initial conditions and a time span
belong to the transient branch; a numerical guess belongs to the stationary
branch. Observation operators connect eligible transient and stationary forward calculations to
MAP estimation and optional sensitivity-based Gauss--Newton, full-Hessian
Laplace or MCMC analysis, with parameter updates
returning to the forward inputs. Stationary calculations use checked equilibria;
implicit sensitivities require a smooth isolated root, while AD-incompatible
closures need explicit finite differences. Forward uncertainty quantification
wraps the same eligible forward problems and observables through
\texttt{UncertainInputs} and \texttt{UQProblem}, without requiring calibration
data. Scenario setup varies physical inputs; Monte Carlo propagation and
non-intrusive polynomial chaos supply output means, covariance, quantiles and
event probabilities, with independent full-model validation of the surrogate.
MC preserves declared joint inputs or posterior draws. PCE uses independent
normal/uniform reference coordinates with an explicit physical-input map;
an arbitrary posterior does not automatically define a PCE basis. This does not
imply universal compatibility of the displayed methods. (b) Required dependencies,
optional extension roles and separate development/figure tools, checked against
the current package declarations. MC and PCE add no optional backend requirement.
The density glyph is conceptual.}
